%HOMEWORK template for M 325K
\documentclass{article}
\headheight=7pt         \topmargin=14pt
\textheight=574pt       \textwidth=445pt
\oddsidemargin=18pt     \evensidemargin=18pt

%Begin package import

\usepackage{amsmath, amssymb, amsthm}
\usepackage{mathtools}
\usepackage{pdfpages}
\usepackage{tikz-cd}

%End package import
%Begin custom commands

\newcommand{\C}{\mathbb{C}}
\newcommand{\R}{\mathbb{R}}
\newcommand{\Q}{\mathbb{Q}}
\newcommand{\Z}{\mathbb{Z}}
\newcommand{\N}{\mathbb{N}}
\newcommand{\abs}[1]{\lvert #1\rvert}
\newcommand{\RM}[1]{\mathrm{#1}}
\newcommand{\BF}[1]{\textbf{#1}}
\newcommand{\e}{\epsilon}
\newcommand{\ceil}[1]{\lceil{#1}\rceil}
\newcommand{\floor}[1]{\lfloor{#1}\rfloor}
\newcommand{\rarr}[0]{\rightarrow}
\newcommand{\IM}[1]{\text{Im}(#1)}
\newcommand{\Hom}[0]{\text{Hom}}
\newcommand{\Pow}[1]{\text{Pow}(#1)}
\newcommand{\angles}[1]{\langle #1 \rangle}

%End custom commands
%Begin custom theorems

\newtheorem{innercustomgeneric}{\customgenericname}
\providecommand{\customgenericname}{}
\newcommand{\newcustomtheorem}[2]{%
\newenvironment{#1}[1]
{%
\renewcommand\customgenericname{#2}%
\renewcommand\theinnercustomgeneric{##1}%
\innercustomgeneric%
}
{\endinnercustomgeneric}
}

\newcustomtheorem{THRM}{Theorem}
\newcustomtheorem{LEM}{Lemma}
\newcustomtheorem{EX}{Exercise}
\newcustomtheorem{COR}{Corollary}
\newcustomtheorem{PROB}{Problem}
\newcustomtheorem{claim}{Claim}

\newenvironment{solution}
{\renewcommand\qedsymbol{$\blacksquare$}\begin{proof}[Solution]}
{\end{proof}}

\newenvironment{definition}
{\renewcommand\qedsymbol{$\blacksquare$}\begin{proof}[Definition]}
{\end{proof}}

%End custom theorems
\begin{document}

\title{Homework 4}
\author{Arham Lodha}%put your name here
\date{February 08, 2024}%enter the due date here
\maketitle

\begin{PROB}{1}\label{Problem 1.}
    Consider $f : A \to B$ a function from the set A to the set B, E subset of B, and F and G subsets of A.
    \begin{enumerate}
        \item Give the mathematical definition of the direct image of F under f.
        \item Give the mathematical definition of the inverse image of E under f.
        \item Show that $$f(E \cup G) = f(E) \cup f(G)$$
    \end{enumerate}
\end{PROB}
\begin{solution}
    \begin{enumerate}
        \item $f(F) = \left\{ f(x) : \forall x \in F \right\}$ 
        \item $f^{-1}(E) = \left\{x \in A : f(x) \in E \right\}$ 
    \end{enumerate}
\end{solution}
\begin{claim}{1c}\label{Claim 1c.}
    Show that $$f(E \cup G) = f(E) \cup f(G)$$
\end{claim}
\begin{proof}
    $\forall b \in f(E \cup G)$, by definition of image, $\exists a \in E \cup G \ni f(a) = b$. Thus by definition of union, $a \in E$ or $a \in G$. If $a \in E$, then $b = f(a) \in f(E)$, by definition of image. If $a \in G$, then $b = f(a) \in f(G)$, by definition of image. Thus $b \in f(E)$ or $b \in f(G)$. Thus by definition of union, $b \in f(E) \cap f(G)$. Thus $f(E \cup G) \subseteq f(E) \cap f(G)$ 
    
    $\forall b \in f(E) \cap f(G)$, by definition of union, $b \in f(E)$ or $b \in f(G)$. If $b \in f(E)$, by definition of image, $\exists a \in E \ni f(a) = b$. If $b \in f(G)$, by definition $\exists a \in G \ni f(a) = b$. Thus $a \in E$ or $a \in G$. Thus $a \in E \cap G$, thus by definition of image, $b \in f(E \cap G)$. Thus $f(E) \cap f(G) \subseteq f(E \cap G)$. Thus $f(E) \cap f(G) = f(E \cap G)$.               
\end{proof}

\bigskip

\begin{PROB}{2a}\label{Problem 2a.}
    Give the definition for a set S to be denumerable
\end{PROB}
\begin{proof}
    S is denumerable if $\exists f: \N \to S$ where $f$ is a bijection.   
\end{proof}
\begin{claim}{2b}\label{Claim 2b.}
    Using only the definition, show that $\Z$ is denumerable.
\end{claim}
\begin{proof}
    $f: \N \to S$ where $\forall n \in \N$, 
    
    $$
    f(n) = \begin{cases}
        -\floor{\frac{n}{2}} & \text{if n is odd} \\
        \frac{n}{2} & \text{if n is even}
    \end{cases}
    $$ 


    $\forall y \in \Z$, if $y \leq 0$, $\exists n \in \N \ni n = -2y + 1$. $n$ is odd thus $f(n) = -\floor{\frac{n}{2}} = -\floor{\frac{-2y + 1}{2}} = -\floor{-y + \frac{1}{2}} = y$. If $y > 0$, then $\exists n \in \N \ni n = 2y$, thus $f(n)  \frac{2y}{2} = y$. Thus $f$ is surjective. 
    
    $\forall a, b \in \N$, suppose $f(a) = f(b)$. If a is odd and b is even, then $f(a) < 0$ and $f(b) > 0$, thus by contradiction both a and b must be both odd or both even. If $a, b$ are odd, then $-\floor{\frac{a}{2}} = -\floor{\frac{b}{2}}$. Thus $\floor{\frac{a}{2}} = \floor{\frac{b}{2}} \implies \frac{a}{2} = \frac{b}{2} \implies a = b$. If $a, b$ are even, then $\frac{a}{2} = f(a) = f(b) = \frac{b}{2}$, thus $a = b$. Thus $f$ is injective.
    
    Thus $f$ is a bijection and $\Z$ is denumerable. 
\end{proof}

\begin{claim}{2c}\label{Claim 2c.}
    Show that $\Q$  is countable 
\end{claim}
\begin{proof}
    Let $f: \Q \to \Z \times \N$ where $\forall \frac{p}{q} \in \Q$, $f(\frac{p}{q}) = (p, q)$. $\forall \frac{p}{q}, \frac{r}{s} \in \Q$, let $f(\frac{p}{q}) = f\left(\frac{r}{s}\right) \implies (p, q) = (r, s) \implies p = r, q =s$. Thus $\frac{p}{q} = \frac{r}{s}$. Thus is injective. 
    You can create a sequence with the elements of $\Z \times \N$ where you are spiraling around the origin, so you start with $(0, 0), (1, 0), (1, 1), (0, 1), (-1, 0), (-1, 0), (-1, -1), (0, -1), (1, -1), (2, -1), (2, 0)\dots$. Since we are going a finite number of elements before we change directions, we don't ever run out of $\N$. Moreover since we are spiraling while making sure to land on lattice points all values in $\Z \times \N$ will be covered. Thus $\exists$ a bijection $\Z \to \N \times \N$ and thus exists inverse bijection $g : \N \times \Z \to \N$.

    $\forall r, s \in \Q \ni (g \circ f)(r) = (g \circ f)(s)$, since $g$ is a bijection it is also injective so $f(r) = f(s)$, since $f$ is injective, $r = s$. Thus $(g \circ f)$ is injective map from $\Q \to \N$. Thus $\Q$ is countable.        
\end{proof}

\bigskip

\begin{PROB}{3}\label{Problem 3.}
    $\forall n \in \N$:
    
    \[
        \sum_{i = 1}^{n} \frac{1}{i(i + 1)} = \frac{1}{1(2)} + \frac{1}{2(3)} + \dots + \frac{1}{n(n + 1)} = \frac{n}{n + 1}
    \] 
\end{PROB}
\begin{proof}
    (Base Case): Let $n = 1$, $\sum_{i = 1}^{1} \frac{1}{i(i + 1)} = \frac{1}{1(2)} = \frac{1}{2}$. $n + 1 = 2$. So $\sum_{i = 1}^{1} \frac{1}{i(i + 1)} = \frac{n}{n+1}$.

    Inductive Hypothesis: Assume $\sum_{i = 1}^{n} \frac{1}{i(i + 1)} = \frac{n}{n + 1}$ is true.

    Inductive Step: $\sum_{i = 1}^{n + 1} \frac{1}{i(i + 1)} = \sum_{i = 1}^{n} \frac{1}{i(i + 1)} + \frac{1}{(n + 1)(n + 2)}$. By Inductive Hypothesis, 
    
    \begin{align*}
        \sum_{i = 1}^{n} \frac{1}{i(i + 1)} + \frac{1}{(n + 1)(n + 2)} &= \frac{n}{n + 1}+ \frac{1}{(n + 1)(n + 2)} \\
        &= \frac{n(n + 2) + 1}{(n + 1)(n+2)} \\
        &= \frac{n^2 + 2n + 1}{(n + 1)(n + 2)} \\
        &= \frac{(n + 1)^2}{(n + 1)(n + 2)} \\
        &= \frac{n + 1}{n + 2}.
    \end{align*}

    Thus, $\sum_{i = 1}^{n + 1} \frac{1}{i(i + 1)} = \frac{n + 1}{n + 2}$. Let $k = n +1$. Thus, $\sum_{i = 1}^{k} \frac{1}{i(i + 1)} = \frac{k}{k + 1}$.    

    Thus by induction, $\forall n \in \N$,
    $\sum_{i = 1}^{n} \frac{1}{i(i + 1)} = \frac{n}{n + 1}$
\end{proof}

\bigskip

\begin{PROB}{4}\label{Problem 4.}
    Using only the nine fundamental algebraic properties of $\R$ , show that:

    \begin{enumerate}
        \item For all $a \in \R$, $-(-a) = a$
        \item For all $a, b \in R$, $(-a)(-b) = ab$ .  
    \end{enumerate}
     
\end{PROB}
\begin{claim}{4a}\label{Claim 4a.}
    For all $a \in \R$, $-(-a) = a$
\end{claim}
\begin{proof}
    $\forall a \in \R$, by the existence of a additive inverse for all reals, $\exists -a \in \R$, $a + (-a) = 0$. By the commutative property of addition, $(-a) + a =0$. Again by the existence of additive inverses $\exists -(-a) \in \R \ni -(-a) + (-a) = 0$. Thus $0 + a = -(-a) + (-a) + a = -(-a)$. By properties of the additive identity, $a = 0 + a = -(-a)$. Thus $a = -(-a)$.         
\end{proof}
\begin{claim}{4b}\label{Claim 4b.}
    For all $a, b \in R$, $(-a)(-b) = ab$ .
\end{claim}
\begin{proof}
$\forall r \in \R$, by the properties of the multiplicative identity, $r * 1 = r$. By the properties of additive identity, $1 + 0 = 1$. Thus $r(1 + 0) = r$. By the distributivity of multiplication over addition, $r*1 + r*0 = r$. By the properties of the additive identity, $r + r*0 = r$. By the existence of a additive inverse, $\exists -r \in \R \ni -r + r = 0$. Thus $0 + r*0 = -r + r + r*0 = -r + r = 0$. By the properties of the additive identity, $0 + r*0 = r*0 =0$. Thus $r*0 = 0$.      

$\forall a, b \in \R$. By the distributivity of multiplication over additon, $b(-1 + 1) = -1(b) + b$. $1 - 1= 0$. Thus $b(0) = -1(b) + b$. By the above, $b(0) = 0$. Thus $-1(b) + b = 0$. By the existence of the additive identity, $-1(b) + b -b = -b \implies -1(b) = -b$. 

Thus $(-a)(-b) = (a)(-1)(b)$. By the commutativity of multiplication, $(-a)(-1)(b) = (-1)(-a)b$. By the above, $(-1)(-a)(b) = -(-a)(b)$. By 4a, $-(-a) = a$. Thus $-(-a)(b) = ab$. Thus $(-a)(-b) = ab$.           
\end{proof}

\end{document}